\numoptstrue
\yearpaper{2026}{Re-Examination \textperiodcentered{} 21 June 2026 \textperiodcentered{} Test Booklet Code 50 \textperiodcentered{} Physics Q 1--45 \textperiodcentered{} 180 Marks}{The NEET (UG) 2026 examination of 3 May 2026 was cancelled by NTA; the result was declared on the re-examination of 21 June 2026, given here.}
\begin{multicols}{2}

\QQ{A particle of mass $M$ moves along a horizontal $x$ axis from $x = 0$ to $x = L$. The coefficient of kinetic friction varies as a function of $x$ as $\mu_k(x) = \mu_0 - \alpha x$, where $\mu_0$, $\alpha$ are constants of appropriate dimensions, so that $\mu_k(L) = 0$. The total work done by the frictional force during the motion is $n\mu_0MgL$, where $g$ is the acceleration due to gravity. The value of $n$ is:\ntag{A}{75}{keph105}{5}{ \textperiodcentered{} work by variable force (area under F-x)}{Class XI Part 1 \textperiodcentered{} Ch 5}}
{\opt{3}{1}{$\dfrac13$}{$\dfrac12$}}

\QQ{The mean free path of molecules in an ideal gas A is half that of another ideal gas B. The diameter of the spherical molecules of gas A is twice the diameter of the molecules of B. If number densities of the gases A and B are $n_A$ and $n_B$, respectively, the correct option is:\ntag{D}{254}{keph205}{11}{ \textperiodcentered{} l=1/(sqrt2 n pi d\textasciicircum{}2)}{Class XI Part 2 \textperiodcentered{} Ch 12}}
{\opt{$n_A = n_B$}{$n_A = 2n_B$}{$n_A = \dfrac14n_B$}{$n_A = \dfrac12n_B$}}

\QQ{Two identical inductors are connected in two different configurations P and Q, where a time varying current $I(t)$ is flowing, as shown in the figure. The induced emf between points $a$ and $b$ for configuration P is $E_P$ and that for configuration Q is $E_Q$. The ratio $E_P/E_Q$ is: [Neglect the effect of mutual inductance.]
\fig{\begin{circuitikz}[scale=0.62,transform shape,font=\Large]
\draw[->] (0,4.6) node[above]{$I(t)$} -- (0,4.1);
\draw (0,4.1) -- (0,3.8) node[circ]{} node[right]{$a$} to[L] (0,2.2) node[circ]{} node[right]{$b$} to[L] (0,0.6) -- (0,0.2);
\node at (0,-0.3) {(P)};
\draw[->] (3.6,4.6) node[above]{$I(t)$} -- (3.6,4.1);
\draw (3.6,4.1) -- (3.6,3.8) node[circ]{} node[right]{$a$};
\draw (2.9,3.8) -- (4.3,3.8); \draw (2.9,3.8) to[L] (2.9,1.6); \draw (4.3,3.8) to[L] (4.3,1.6);
\draw (2.9,1.6) -- (4.3,1.6); \draw (3.6,1.6) node[circ]{} node[below right]{$b$} -- (3.6,1.0);
\node at (3.6,-0.3) {(Q)};
\end{circuitikz}}\ntag{A}{168}{leph106}{15}{}{Class XII Part 1 \textperiodcentered{} Ch 6}}
{\opt{$\dfrac14$}{$\dfrac12$}{1}{2}}

\QQ{For sound waves, if the number of nodes for the $5^{\text{th}}$ harmonic of an open-ended pipe is $n$ and that for the $9^{\text{th}}$ harmonic of the same pipe with one of its ends closed is $m$, the ratio $\dfrac nm$ is\ntag{A}{292}{keph207}{15}{}{Class XI Part 2 \textperiodcentered{} Ch 14}}
{\opt{$\dfrac59$}{$\dfrac95$}{1}{$\dfrac35$}}

\QQ{Consider a long solenoid of length $l$ and radius $r$. If $n$ is the number of turns per unit length and $\mu_0$ is the permeability of free space, the inductance of the solenoid is:\ntag{D}{168}{leph106}{15}{ \textperiodcentered{} L=mu0 n\textasciicircum{}2 A l}{Class XII Part 1 \textperiodcentered{} Ch 6}}
{\opt{$\mu_0\pi n^2r^2l$}{$\mu_0n^2r^2l$}{$(\mu_0/2\pi)n^2r^2l$}{$2\mu_0\pi n^2r^2l$}}

\QQ{Consider a particle moving along a straight line, whose position as a function of time is given by $s(t) = \alpha t^2-\beta t+\gamma$, where $\alpha = 1\unit{m\,s^{-2}}$, $\beta = 6\unit{m\,s^{-1}}$ and $\gamma = 5\unit{m}$. The average speed of the particle, in $\mathrm{m\,s^{-1}}$ from $t = 0$ to $t = 6\unit{s}$ is:\ntag{A}{21}{keph102}{9}{}{Class XI Part 1 \textperiodcentered{} Ch 2}}
{\opt{12}{6}{3}{0}}

\QQ{Consider the following nuclear reaction
\[{}^{238}\mathrm U\longrightarrow{}^{234}\mathrm{Th}+{}^{4}\mathrm{He}\]
Take masses of ${}^{238}$U, ${}^{234}$Th and ${}^{4}$He as $238.050\unit{u}$, $234.043\unit{u}$ and $4.003\unit{u}$, respectively. The Q value for the reaction, in keV, is: [Given: $1\unit{u} = 931.5\unit{MeV\,c^{-2}}$]\ntag{D}{319}{leph205}{14}{ \textperiodcentered{} Q=(sum mi - sum mf)c\textasciicircum{}2}{Class XII Part 2 \textperiodcentered{} Ch 13}}
{\opt{3726}{3730}{3736}{3740}}

\QQ{A beam of light falls on a metal surface such that photo-electrons are generated. If power of the light source starts to decrease linearly with time $t$, then variation of the photocurrent $I$ and magnitude of the stopping potential $|V|$ with time is best represented by:\ntag{D}{278}{leph203}{5}{ \textperiodcentered{} Figs 11.2, 11.3}{Class XII Part 2 \textperiodcentered{} Ch 11}}
{\optfig{\gaxes{$t$}{$I$}{\draw (0,1.5) -- (1.6,0.1);}\,\gaxes{$t$}{$|V|$}{\draw (0,1.2) -- (1.8,1.2);}}
{\gaxes{$t$}{$I$}{\draw (0,1.5) -- (1.6,0.1);}\,\gaxes{$t$}{$|V|$}{\draw (0,1.5) -- (1.6,0.1);}}
{\gaxes{$t$}{$I$}{\draw (0,1.2) -- (1.8,1.2);}\,\gaxes{$t$}{$|V|$}{\draw (0,1.5) -- (1.6,0.1);}}
{\gaxes{$t$}{$I$}{\draw (0,1.2) -- (1.8,1.2);}\,\gaxes{$t$}{$|V|$}{\draw (0,1.2) -- (1.8,1.2);}}}

\QQ{Three identical capacitors, P, Q and S, each of the capacitance $C$, are connected to a battery of voltage $V$, as shown in the figure. If the energy stored in the capacitor P and total energy stored in the system are $U_P$ and $U_T$, respectively, then the ratio $\dfrac{U_P}{U_T}$ is:
\fig{\begin{circuitikz}[scale=0.62,transform shape,font=\Large]
\draw (0,2.4) to[C,l=P] (2.2,2.4) to[C,l=Q] (4.4,2.4) -- (4.4,1.2);
\draw (0,2.4) -- (0,1.2) to[C,l_=S] (4.4,1.2);
\draw (0,1.2) -- (-0.6,1.2) -- (-0.6,0) to[battery1,l_=$V$] (5.0,0) -- (5.0,1.2) -- (4.4,1.2);
\end{circuitikz}}\ntag{A}{72}{leph102}{28}{}{Class XII Part 1 \textperiodcentered{} Ch 2}}
{\opt{2/3}{1/3}{1/2}{1/6}}

\QQ{Water flows in a streamline motion through a horizontal pipe of circular cross-section as shown in the figure. The pressure difference of water between P and Q is $15\unit{N\,m^{-2}}$. The area of cross-section at P and Q are $40\unit{cm^2}$ and $20\unit{cm^2}$, respectively. The rate of flow of water through the pipe, in $\mathrm{cm^3\,s^{-1}}$, is: [Take density of water $= 1000\unit{kg\,m^{-3}}$]
\fig{\begin{tikzpicture}[scale=0.8]
\draw (0,0.6) -- (1.4,0.6) .. controls (2.0,0.6) and (2.2,0.3) .. (2.8,0.3) .. controls (3.4,0.3) and (3.6,0.6) .. (4.2,0.6);
\draw (0,-0.6) -- (1.4,-0.6) .. controls (2.0,-0.6) and (2.2,-0.3) .. (2.8,-0.3) .. controls (3.4,-0.3) and (3.6,-0.6) .. (4.2,-0.6);
\node[lbl] at (0.8,0) {P}; \node[lbl] at (2.8,0) {Q};
\end{tikzpicture}}\ntag{A}{188}{keph202}{9}{}{Class XI Part 2 \textperiodcentered{} Ch 9}}
{\opt{100}{200}{300}{400}}

\QQ{A current $I_0$ flows through a metallic circular loop of radius $r$ as shown in the figure. Resistance of the segment ABC is half that of ADC. Magnitude of magnetic field at the centre O of the loop is:
\fig{\begin{tikzpicture}[scale=0.8]
\draw (0,0) circle (1.1);
\draw[->] (0,2.0) -- (0,1.4) node[pos=0,right,lbl]{$I_0$}; \draw (0,1.4) -- (0,1.1);
\draw[->] (0,-1.1) -- (0,-1.8) node[right,lbl]{$I_0$};
\draw[densely dashed] (0,1.1) -- (0,-1.1); \draw[densely dashed,->] (0,0) -- (-40:1.1) node[pos=0.6,right,lbl]{$r$};
\fill (0,0) circle (1.3pt) node[left,lbl]{O};
\node[lbl,above left] at (0,1.1) {A}; \node[lbl,below left] at (0,-1.1) {C};
\node[lbl,left] at (-1.1,0) {B}; \node[lbl,right] at (1.1,0) {D};
\end{tikzpicture}}\ntag{A}{116}{leph104}{10}{}{Class XII Part 1 \textperiodcentered{} Ch 4}}
{\opt{$\dfrac{\mu_0I_0}{12r}$}{$\dfrac{\mu_0I_0}{4r}$}{$\dfrac{\mu_0I_0}{2r}$}{$\dfrac{\mu_0I_0}{2\pi r}$}}

\QQ{In the measurement of viscosity of liquids using terminal velocity experiment, spherical balls of same radius but having different densities are used. The variation of the terminal velocity ($v$) with the ratio of density of spherical ball ($\sigma$) to density of the liquid ($\rho$), is best represented by:\ntag{D}{192}{keph202}{13}{ \textperiodcentered{} vt=2a\textasciicircum{}2(rho-sigma)g/9eta}{Class XI Part 2 \textperiodcentered{} Ch 9}}
{\optfig{\gaxes{$\sigma/\rho$}{$v$}{\draw (0.7,0) -- (1.8,1.2); \node[font=\tiny,below] at (0.7,0) {1}; \node[font=\tiny,below] at (0,0) {0};}}
{\gaxes{$\sigma/\rho$}{$v$}{\draw[densely dashed] (0,0) -- (0.7,0.45); \draw (0.7,0.45) -- (1.8,1.15); \node[font=\tiny,below] at (0.7,0) {1}; \node[font=\tiny,below] at (0,0) {0};}}
{\gaxes{$\sigma/\rho$}{$v$}{\draw[densely dashed] (0.3,0) -- (0.7,0.35); \draw (0.7,0.35) -- (1.8,1.3); \node[font=\tiny,below] at (0.7,0) {1}; \node[font=\tiny,below] at (0,0) {0};}}
{\gaxes{$\sigma/\rho$}{$v$}{\draw (0.7,0) -- (0.7,1.0) -- (1.8,1.0); \node[font=\tiny,below] at (0.7,0) {1}; \node[font=\tiny,below] at (0,0) {0};}}}

\QQ{Two planets $P_1$ and $P_2$ with equal mass have radii $R_1$ and $R_2$, respectively, where $R_2 = \dfrac{R_1}{2}$. The escape speeds of $P_1$ and $P_2$ are $v_1$ and $v_2$, respectively. Then $\dfrac{v_2}{v_1}$ is:\ntag{D}{136}{keph107}{10}{ \textperiodcentered{} v=sqrt(2GM/R)}{Class XI Part 1 \textperiodcentered{} Ch 7}}
{\opt{$\dfrac1{\sqrt2}$}{1}{$\sqrt2$}{2}}

\QQ{In a solar system, the time-period of revolution of a planet tracing a circular orbit of radius $R$ is proportional to:\ntag{D}{129}{keph107}{3}{ \textperiodcentered{} T\textasciicircum{}2 proportional to a\textasciicircum{}3}{Class XI Part 1 \textperiodcentered{} Ch 7}}
{\opt{$R^{1/2}$}{$R^{3/2}$}{$R^2$}{$R^3$}}

\QQ{Two infinitely long parallel conducting wires A and B carry currents $I$ and $2I$, respectively, in the same direction. The wire A has uniform mass per unit length $\lambda$ and lies on an insulated floor. The wire B is kept fixed at a height $h$ above the floor. The minimum magnitude of $h$ so that the wire A does not rise from the floor is: [$g$ is the acceleration due to gravity and $\mu_0$ is the permeability of free space.]\ntag{A}{123}{leph104}{17}{}{Class XII Part 1 \textperiodcentered{} Ch 4}}
{\opt{$\dfrac{\mu_0I^2}{2\pi\lambda g}$}{$\dfrac{\mu_0I^2}{\pi\lambda g}$}{$\dfrac{2\mu_0I^2}{\pi\lambda g}$}{$\dfrac{4\mu_0I^2}{\pi\lambda g}$}}

\QQ{An ideal Zener diode with breakdown voltage of $-3\unit{V}$ is reverse biased with a negative input voltage $V_i = -5\unit{V}$. The magnitude of voltage difference between point B and A is:
\fig{\begin{circuitikz}[scale=0.62,transform shape,font=\Large]
\draw (0,2) node[circ]{} node[above]{C} to[zDo] (1.6,2) node[circ]{} node[above]{B} to[R,l=R] (3.6,2) node[circ]{} node[above]{A} -- (4.2,2) -- (4.2,0) to[battery1] (0,0) -- (0,2);
\end{circuitikz}}\ntag{X}{}{}{}{}{Zener diode}}
{\opt{3\,V}{2\,V}{1\,V}{0\,V}}

\QQ{In an adiabatic expansion, the temperature of one mole of an ideal monatomic gas ($\gamma = 5/3$) decreases from $60\unit{K}$ to $50\unit{K}$. The work done by the gas in the process is: (Take the universal gas constant as $R = 8.3\unit{J\,mol^{-1}K^{-1}}$)\ntag{D}{235}{keph204}{10}{ \textperiodcentered{} W=muR(T1-T2)/(gamma-1)}{Class XI Part 2 \textperiodcentered{} Ch 11}}
{\opt{41.5\,J}{83\,J}{124.5\,J}{166\,J}}

\QQ{A ray of light with wavelength $\lambda$ is incident on three different photoelectric cells namely 1, 2 and 3. The threshold wavelength of these photoelectric cells are $\lambda_1$, $\lambda_2$, and $\lambda_3$, respectively and the magnitude of stopping potentials of these cells are $V_1$, $V_2$ and $V_3$, respectively. The relation between $\lambda$ and threshold wavelengths are $\lambda_1<\lambda$, $\lambda_2>\lambda$ and $\lambda_3\gg\lambda$. The correct option is:\ntag{A}{281}{leph203}{8}{}{Class XII Part 2 \textperiodcentered{} Ch 11}}
{\opt{$V_1 = 0$, $V_2<V_3$}{$V_1 = 0$, $V_2>V_3$}{$V_1>V_2$, $V_3 = 0$}{$V_1<V_2$, $V_3 = 0$}}

\QQ{A photon and an electron, each of $20\unit{eV}$ energy, move in free space. The ratio of linear momentum of electron $p_e$ to that of photon $p_{Ph}$, $\dfrac{p_e}{p_{Ph}}$ is: [Take speed of light $= 3\times10^8\unit{m\,s^{-1}}$, charge of electron $= -1.6\times10^{-19}\unit{C}$ and mass of electron $= 9\times10^{-31}\unit{kg}$]\ntag{A}{283}{leph203}{10}{}{Class XII Part 2 \textperiodcentered{} Ch 11}}
{\opt{$\dfrac2{450}$}{$\dfrac1{250}$}{225}{275}}

\QQ{Which of the following measurements require `index correction'?\ntag{A}{6}{keph101}{6}{ \textperiodcentered{} Lab manual: index correction on the optical bench}{Class XI Part 1 \textperiodcentered{} Ch 1}}
{\opt{Measurement of resistance of a wire using meter bridge}{Measurement of gravitational acceleration using simple pendulum}{Measurement of focal length of lenses using optical bench}{Measurement of speed of sound using resonance tube}}

\QQ{A unit positive point charge is taken slowly through an infinitesimally thin tube that is inside a charged dielectric sphere of radius $R$, having uniform positive charge density $\rho$, as shown in the figure. The initial and final positions of the charge are marked by A and B at distance $2R$ and $3R$ respectively, from the centre of the sphere. In this process, the magnitude of the total work done on the point charge is $\dfrac{\rho R^2}{n\varepsilon_0}$. The value of $n$ is: ($\varepsilon_0$ is the permittivity of vacuum)
\fig{\begin{tikzpicture}[scale=0.75]
\draw (0,0) circle (1);
\fill (0,0) circle (1.2pt) node[below,lbl]{O}; \node[lbl] at (-0.35,-0.25) {$\rho$};
\coordinate (A) at (1.8,0.85); \coordinate (B) at (-2.95,0.55);
\fill (A) circle (1.5pt) node[right,lbl]{A}; \fill (B) circle (1.5pt) node[left,lbl]{B};
\draw[densely dashed,->] (A) -- (-1.0,0.68); \draw[densely dashed] (-1.0,0.68) -- (B);
\draw (0,0) -- (A) node[pos=0.6,above,lbl]{$2R$}; \draw (0,0) -- (B) node[pos=0.55,below,lbl]{$3R$};
\end{tikzpicture}}\ntag{A}{52}{leph102}{8}{}{Class XII Part 1 \textperiodcentered{} Ch 2}}
{\opt{2}{6}{9}{18}}

\QQ{Consider three media P, Q and R with refractive indices 1, 1.25, and 1.5 respectively. The medium Q having a thickness of $5\unit{cm}$ is placed between extended media P and R as shown in the figure. An object O is placed at the centre of medium Q. If viewed from medium P near the normal direction, the apparent depth of O is $h_1$. For similar observation from medium R, the apparent depth is $h_2$. The value of $|h_1-h_2|$, in cm, is:
\fig{\begin{tikzpicture}[scale=0.75]
\fill[gray!35] (0,0) rectangle (3,1.2);
\draw (0,0) rectangle (3,2.4); \draw[densely dotted] (0,2.4) -- (0,4) (3,2.4) -- (3,4);
\draw (0,1.2) -- (3,1.2);
\node[lbl] at (2.3,3.0) {$n_P = 1$}; \node[lbl] at (2.3,1.55) {$n_Q = 1.25$}; \node[lbl] at (2.3,0.35) {$n_R = 1.5$};
\draw[<->] (-0.3,1.2) -- (-0.3,2.4) node[midway,left,lbl]{5\,cm};
\fill (1.4,1.8) circle (1.5pt) node[above,lbl]{O};
\end{tikzpicture}}\ntag{A}{229}{leph201}{9}{}{Class XII Part 2 \textperiodcentered{} Ch 9}}
{\opt{0}{1}{2}{3}}

\QQ{A frictionless circular wire of unit radius is fixed on the horizontal plane. Two point particles of unit mass start moving simultaneously from point A ($\theta = \frac\pi2$) with identical uniform angular speeds in opposite directions, and meet again at point B ($\theta = -\frac\pi2$). During this time, which of the following figures schematically represent the magnitude of the total linear momentum $P$ of the system, as a function of $\theta$?
\fig{\begin{tikzpicture}[scale=0.7]
\draw (0,0) circle (1); \draw[densely dashed] (-1.5,0) -- (1.5,0) (0,-1.3) -- (0,1.3);
\node[lbl,above left] at (0,1) {A}; \node[lbl,below left] at (0,-1) {B};
\draw[densely dashed] (0,0) -- (45:1); \draw (0,0.4) arc (90:45:0.4); \node[lbl] at (0.35,0.55) {$\theta$};
\draw[->] (100:1.15) arc (100:130:1.15); \draw[->] (80:1.15) arc (80:50:1.15);
\end{tikzpicture}}\ntag{A}{53}{keph104}{5}{}{Class XI Part 1 \textperiodcentered{} Ch 4}}
{\optfig{\gaxes{$\theta$}{$P$}{\draw (0,1.2) -- (1.6,1.2); \node[font=\tiny,below] at (0,0) {$\frac\pi2$}; \node[font=\tiny,below] at (1.6,0) {$-\frac\pi2$};}}
{\gaxes{$\theta$}{$P$}{\draw (0,1.3) .. controls (0.2,0.4) and (0.6,0.05) .. (0.8,0) .. controls (1.0,0.05) and (1.4,0.4) .. (1.6,1.3); \node[font=\tiny,below] at (0,0) {$\frac\pi2$}; \node[font=\tiny,below] at (1.6,0) {$-\frac\pi2$};}}
{\gaxes{$\theta$}{$P$}{\draw (0,0) .. controls (0,0.9) and (0.35,1.3) .. (0.8,1.3) .. controls (1.25,1.3) and (1.6,0.9) .. (1.6,0); \node[font=\tiny,below] at (0,0) {$\frac\pi2$}; \node[font=\tiny,below] at (1.6,0) {$-\frac\pi2$};}}
{\gaxes{$\theta$}{$P$}{\draw (0,1.4) -- (1.6,0); \node[font=\tiny,below] at (0,0) {$\frac\pi2$}; \node[font=\tiny,below] at (1.6,0) {$-\frac\pi2$};}}}

\QQ{The temperature of a metallic sphere of radius $R$ is increased by a small amount $\Delta T$. If the linear coefficient of thermal expansion of the metal is $\alpha$, the approximate increase in the volume of the sphere is:\ntag{D}{207}{keph203}{6}{ \textperiodcentered{} alpha\_V=3 alpha\_l}{Class XI Part 2 \textperiodcentered{} Ch 10}}
{\opt{$2\pi R^3\alpha\Delta T$}{$3\pi R^3\alpha\Delta T$}{$4\pi R^3\alpha\Delta T$}{$6\pi R^3\alpha\Delta T$}}

\QQ{A cylindrical cork of uniform density floats in a liquid of density $\rho_1$. If the cork is depressed slightly and released, it oscillates harmonically with time period $T$. If the same cork floats in another liquid of density $\rho_2$, then the similar oscillation has time period $2T$. The value of $\rho_2/\rho_1$ is:\ntag{D}{277}{keph206}{19}{ \textperiodcentered{} Exercise 13.17}{Class XI Part 2 \textperiodcentered{} Ch 13}}
{\opt{4}{2}{1/2}{1/4}}

\QQ{One main scale division of a Vernier calliper is equal to $1\unit{mm}$ and the number of divisions on the Vernier scale is 10. When both the jaws touch each other, the Vernier scale shifts to the left of zero of the main scale in such a way that $4^{\text{th}}$ Vernier division coincides with a division of the main scale. If this Vernier calliper measures the length of the wire to be $1\unit{cm}$, the actual length of the wire is:\ntag{A}{11}{keph101}{11}{ \textperiodcentered{} Lab manual: vernier zero error}{Class XI Part 1 \textperiodcentered{} Ch 1}}
{\opt{0.60\,cm}{0.96\,cm}{1.00\,cm}{1.04\,cm}}

\QQ{A solid sphere A of radius $R$ and mass $M$ is attached at a point to a smaller solid sphere B of radius $r<R$ and mass $m<M$. Assume that the line joining their centres lies along the horizontal. The moment of inertia of the system calculated about a vertical axis passing through the centre of A is $I_A$ and that calculated about a vertical axis passing through the centre of B is $I_B$. The difference $I_A-I_B$ is:
\fig{\begin{tikzpicture}[scale=0.65]
\draw (0,0) circle (1); \draw (1.35,0) circle (0.35); \draw[densely dashed] (0,-1.5) -- (0,1.5) node[right,lbl]{$I_A$}; \draw[densely dotted] (0,0) -- (1.35,0);
\node[lbl] at (-1.15,0.75) {A}; \node[lbl] at (1.8,0.45) {B};
\begin{scope}[xshift=4.2cm]
\draw (0,0) circle (1); \draw (1.35,0) circle (0.35); \draw[densely dashed] (1.35,-1.5) -- (1.35,1.5) node[right,lbl]{$I_B$}; \draw[densely dotted] (0,0) -- (1.35,0);
\node[lbl] at (-1.15,0.75) {A}; \node[lbl] at (1.8,0.45) {B};
\end{scope}
\end{tikzpicture}}\ntag{X}{}{}{}{}{parallel axis theorem}}
{\opt{$(M-m)(R+r)^2$}{$(m-M)(R+r)^2$}{$(m-M)(R-r)^2$}{0}}

\QQ{Consider a spring-mass simple harmonic oscillator in one dimension. The mass of the particle is $m\unit{kg}$ and the spring constant is $k\unit{N\,m^{-1}}$. At a given instant, the extension of the spring is $x$ metre and the speed of the particle is $v\unit{m\,s^{-1}}$. On the $x$--$v$ plane, if the graph of $v$ as a function of $x$ is a circle, then the correct option is:\ntag{A}{266}{keph206}{8}{}{Class XI Part 2 \textperiodcentered{} Ch 13}}
{\opt{$k = \dfrac1m$}{$k = m$}{$k = m^2$}{$k = \sqrt m$}}

\QQ{The lens combination as shown in the figure, consists of two lenses, $L_1$ and $L_2$, of the focal lengths $+10\unit{cm}$ and $-10\unit{cm}$, respectively. The position of the image formed is:
\fig{\begin{tikzpicture}[scale=0.75]
\draw[densely dotted] (-0.4,0) -- (4.6,0);
\draw[->,thick] (0,0) -- (0,0.8) node[left,lbl]{} ; \node[lbl,below] at (0,0) {O};
\draw (2.6,0) ellipse (0.18 and 1.2); \node[lbl] at (2.6,1.45) {$L_1$};
\draw (3.05,1.2) -- (3.35,1.2) .. controls (3.15,0.4) and (3.15,-0.4) .. (3.35,-1.2) -- (3.05,-1.2) .. controls (3.25,-0.4) and (3.25,0.4) .. (3.05,1.2); \node[lbl] at (3.2,1.45) {$L_2$};
\draw[<->] (0,-1.5) -- (2.6,-1.5) node[midway,below,lbl]{30\,cm}; \draw[<->] (2.6,-1.5) -- (3.2,-1.5) node[midway,below,lbl]{3\,cm};
\end{tikzpicture}}\ntag{A}{238}{leph201}{18}{}{Class XII Part 2 \textperiodcentered{} Ch 9}}
{\opt{20\,cm to the left of the concave lens}{60\,cm to the left of the concave lens}{30\,cm to the right of the concave lens}{60\,cm to the right of the concave lens}}

\QQ{An ac voltage $V = 220\sin(2\times10^3t)$ volt is applied to a series LCR circuit. Then the current amplitude in this circuit is: (Given: $L = 10\unit{mH}$, $C = 25\unit{\mu F}$, $R = 100\unit{\Omega}$)\ntag{A}{188}{leph107}{12}{}{Class XII Part 1 \textperiodcentered{} Ch 7}}
{\opt{2.2\,A}{5.5\,A}{11.0\,A}{22.0\,A}}

\QQ{A thin horizontal disc is rotating about a vertical axis passing through its fixed centre O. Its angular momentum is $L_A$ and $L_B$ computed about points A and B, respectively, with $OB = 2\times OA$. The value of $\dfrac{L_A}{L_B}$ is:
\fig{\begin{tikzpicture}[scale=0.7]
\draw (0,0) circle (1); \draw[->] (130:1.15) arc (130:160:1.15);
\fill (0,0) circle (1.3pt) node[below,lbl]{O}; \fill (1.1,0) circle (1.3pt) node[below,lbl]{A}; \fill (2.2,0) circle (1.3pt) node[below,lbl]{B};
\draw[densely dotted] (0,0) -- (2.2,0);
\end{tikzpicture}}\ntag{A}{108}{keph106}{17}{}{Class XI Part 1 \textperiodcentered{} Ch 6}}
{\opt{$\dfrac14$}{$\dfrac12$}{1}{2}}

\QQ{A conducting loop of finite resistance lies on the $x$--$y$ plane. There is a constant magnetic field in the $z$ direction. The area of the loop varies with time $t$, as $A = A_0(1+\sin t)$ in appropriate units. The figure that correctly indicates the qualitative behaviour of the power $P$ dissipated in the loop as a function of time is:\ntag{A}{157}{leph106}{4}{}{Class XII Part 1 \textperiodcentered{} Ch 6}}
{\optfig{\gaxes{$t$}{$P$}{\draw (0,0) .. controls (0.1,1.4) and (0.7,1.4) .. (0.8,0) .. controls (0.9,1.4) and (1.5,1.4) .. (1.6,0); \node[font=\tiny,below] at (0,0) {0};}}
{\gaxes{$t$}{$P$}{\draw (0,1.3) .. controls (0.2,0.4) and (0.6,0.05) .. (0.8,0) .. controls (1.0,0.05) and (1.4,0.4) .. (1.6,1.3); \node[font=\tiny,below] at (0,0) {0};}}
{\gaxes{$t$}{$P$}{\draw (0,1.2) -- (1.8,1.2); \node[font=\tiny,below] at (0,0) {0};}}
{\gaxes{$t$}{$P$}{\draw (0,0) .. controls (0.9,0.1) and (1.4,0.5) .. (1.8,1.5); \node[font=\tiny,below] at (0,0) {0};}}}

\QQ{A point charge $Q$ is placed inside a cavity within a solid isolated conducting sphere. Consider points A, B and C as shown in the figure, where the magnitudes of the electric fields are $E_A$, $E_B$, $E_C$, respectively. The points B and C are at the same distance from the center of the solid sphere. The correct option is:
\fig{\begin{tikzpicture}[scale=0.75]
\fill[gray!35] (0,0) circle (1.2); \draw (0,0) circle (1.2);
\fill[white] (0.45,0.1) circle (0.4); \draw (0.45,0.1) circle (0.4);
\fill (0.45,0.25) circle (1.2pt) node[above,lbl]{$Q$}; \fill (0.55,-0.1) circle (1pt) node[right,lbl]{A};
\draw[densely dotted] (-2.0,0) -- (2.0,0);
\fill (-1.7,0) circle (1.3pt) node[above,lbl]{C}; \fill (1.7,0) circle (1.3pt) node[above,lbl]{B};
\end{tikzpicture}}\ntag{A}{63}{leph102}{19}{}{Class XII Part 1 \textperiodcentered{} Ch 2}}
{\opt{$E_A = 0$, $E_B = E_C$}{$E_A\ne0$, $E_B = E_C$}{$E_A = 0$, $E_B>E_C$}{$E_A\ne0$, $E_B<E_C$}}

\QQ{Consider a fixed uniformly charged insulating sphere with radius $R$ and total charge $+Q$. A point charge $-q$ ($q\ll Q$) with mass $m$ is released from rest at a distance of $3R$ from the centre of the charged sphere. When the point charge reaches the surface of the sphere, its speed is: ($\varepsilon_0$ is the permittivity of vacuum, neglect gravitational forces).\ntag{A}{55}{leph102}{11}{}{Class XII Part 1 \textperiodcentered{} Ch 2}}
{\opt{$\sqrt{\dfrac{3Qq}{4\pi\varepsilon_0mR}}$}{$\sqrt{\dfrac{2Qq}{3\pi\varepsilon_0mR}}$}{$\sqrt{\dfrac{Qq}{3\pi\varepsilon_0mR}}$}{$\sqrt{\dfrac{Qq}{4\pi\varepsilon_0mR}}$}}

\QQ{In Geiger--Marsden experiment, the number of scattered $\alpha$-particles $N(\theta)$ is plotted as a function of scattering angle $\theta$. Which of the following options represents the correct plot? ($N(\theta)$ on a logarithmic scale, $\theta$ from $0^\circ$ to $180^\circ$.)\ntag{D}{293}{leph204}{4}{ \textperiodcentered{} Fig 12.3}{Class XII Part 2 \textperiodcentered{} Ch 12}}
{\optfig{\gaxes{$\theta$}{$N(\theta)$}{\draw (0.15,0.3) -- (1.8,1.5);}}
{\gaxes{$\theta$}{$N(\theta)$}{\draw (0.05,1.5) .. controls (0.8,1.55) and (1.4,1.0) .. (1.8,0.3);}}
{\gaxes{$\theta$}{$N(\theta)$}{\draw (0.05,1.6) .. controls (0.3,0.6) and (0.8,0.3) .. (1.8,0.2);}}
{\gaxes{$\theta$}{$N(\theta)$}{\draw (0.1,0.1) .. controls (0.5,0.2) and (0.5,1.5) .. (0.95,1.5) .. controls (1.4,1.5) and (1.4,0.2) .. (1.8,0.1);}}}

\QQ{Consider two circuits, (A) and (B), each having two resistors. One of them has a positive temperature coefficient of resistance, $+\alpha$, while the other one has a negative temperature coefficient, $-\alpha$, as shown in the figure. The current through these circuits are denoted by $I_A$ and $I_B$. At initial temperature, the resistance of the two resistors is $R_0$. As the temperature is increased, the correct option that describes the variation of current in these circuits is:
\fig{\begin{circuitikz}[scale=0.55,transform shape,font=\Large]
\draw (0,1.6) to[battery1,l=$V$] (2.2,1.6) to[short,i=$I_A$] (4.4,1.6) -- (4.4,0) to[R,l={$R_0(1+\alpha\Delta T)$}] (2.2,0) to[R,l={$R_0(1-\alpha\Delta T)$}] (0,0) -- (0,1.6);
\node at (2.2,-1.6) {(A)};
\draw (6.0,2.6) to[battery1,l=$V$] (8.0,2.6) to[short,i=$I_B$] (9.4,2.6) -- (9.4,0);
\draw (6.0,2.6) -- (6.0,0); \draw (6.0,1.3) to[R,l={$R_0(1-\alpha\Delta T)$}] (9.4,1.3); \draw (6.0,0) to[R,l_={$R_0(1+\alpha\Delta T)$}] (9.4,0);
\node at (7.7,-1.6) {(B)};
\end{circuitikz}}\ntag{A}{90}{leph103}{10}{}{Class XII Part 1 \textperiodcentered{} Ch 3}}
{\opt{$I_A$ remains constant while $I_B$ increases}{$I_A$ decreases while $I_B$ increases}{$I_A$ increases while $I_B$ decreases}{Both $I_A$ and $I_B$ remain constant}}

\QQ{Consider that $\sigma_s$, $k_B$, $b$ represents Stefan--Boltzmann constant, Boltzmann constant and Wien's displacement law constant, respectively. The dimension of $\sigma_sk_B^{-1}b$ is\ntag{A}{7}{keph101}{7}{}{Class XI Part 1 \textperiodcentered{} Ch 1}}
{\opt{$[\mathrm{L^{-1}T^{-1}K^{-2}}]$}{$[\mathrm{L^{-1}K^{-2}}]$}{$[\mathrm{L^{-1}T^{-1}K^{-3}}]$}{$[\mathrm{L^{-1}T^{-1}K^{-4}}]$}}

\QQ{An electromagnetic wave travelling in a lossless dielectric medium having a dielectric constant $\varepsilon_r = 9$, has the electric field $E_x = E_0\sin(kz-2\pi\times10^6t)\unit{V\,m^{-1}}$ where $E_0$ is the amplitude and $k$ is the wave vector. Among the following options, the \textbf{incorrect} choice is\ntag{A}{207}{leph108}{7}{}{Class XII Part 1 \textperiodcentered{} Ch 8}}
{\opt{The speed of the electromagnetic wave inside the medium is $10^8\unit{m\,s^{-1}}$}{The wavelength of the electromagnetic wave inside the medium is $300\unit{m}$}{The magnetic field is given by the relation $B_y = \dfrac{B_0}{v}\sin(kz-2\pi\times10^6t)$ where $v$ is the speed of the electromagnetic wave inside the medium}{The direction of propagation of the electromagnetic wave is along $+z$}}

\QQ{One mole of an ideal monatomic gas undergoes a cyclic process as shown in the figure. The total heat supplied to the gas is:
\fig{\begin{tikzpicture}[xscale=0.5,yscale=0.008]
\draw[->] (0,0) -- (7,0) node[right,lbl]{$V\,(\mathrm{m^3})$}; \draw[->] (0,0) -- (0,380) node[above,lbl]{$P\,(\mathrm{N\,m^{-2}})$};
\draw[densely dashed] (0,300) node[left,lbl]{300} -- (2,300) (0,100) node[left,lbl]{100} -- (2,100) (2,100) -- (2,0) node[below,lbl]{2} (5,100) -- (5,0) node[below,lbl]{5};
\draw[thick,postaction={decorate},decoration={markings,mark=at position 0.125 with {\arrow{Stealth}},mark=at position 0.375 with {\arrow{Stealth}},mark=at position 0.625 with {\arrow{Stealth}},mark=at position 0.875 with {\arrow{Stealth}}}] (2,300) -- (5,300) -- (5,100) -- (2,100) -- cycle;
\node[lbl,above left] at (2,300) {$a$}; \node[lbl,above right] at (5,300) {$b$}; \node[lbl,below right] at (5,100) {$c$}; \node[lbl,below left] at (2,100) {$d$};
\end{tikzpicture}}\ntag{D}{236}{keph204}{11}{ \textperiodcentered{} cyclic process: Q=W=area}{Class XI Part 2 \textperiodcentered{} Ch 11}}
{\opt{400\,J}{500\,J}{600\,J}{800\,J}}

\QQ{Consider that an electron is revolving in an excited state of hydrogen atom with velocity $\sqrt{25.6}\times10^5\unit{m\,s^{-1}}$. The radius of the orbit is $x\times10^{-9}\unit{m}$. The value of $x$ is: [Take the mass of electron to be $9\times10^{-31}\unit{kg}$, charge of electron $= -1.6\times10^{-19}\unit{C}$ and $\dfrac1{4\pi\varepsilon_0} = 9\times10^9\unit{N\,m^2C^{-2}}$]\ntag{D}{295}{leph204}{6}{ \textperiodcentered{} r=e\textasciicircum{}2/(4 pi eps0 m v\textasciicircum{}2)}{Class XII Part 2 \textperiodcentered{} Ch 12}}
{\opt{4}{3}{2}{1}}

\QQ{A car travels on a circular racetrack of radius $50\unit{m}$, which is banked at an angle $\theta$. If the car travels at a speed $10\unit{m\,s^{-1}}$, then the wear and tear on its tyres is minimum. Taking the acceleration due to gravity to be $10\unit{m\,s^{-2}}$, the value of $\theta$ is:\ntag{D}{64}{keph104}{16}{ \textperiodcentered{} v0=sqrt(Rg tan theta)}{Class XI Part 1 \textperiodcentered{} Ch 4}}
{\opt{$\tan^{-1}\!\left(\dfrac15\right)$}{$\tan^{-1}\!\left(\dfrac25\right)$}{$\tan^{-1}(\sqrt3/2)$}{$\tan^{-1}(2\sqrt3)$}}

\QQ{Three identical p-n junction diodes $D_1$, $D_2$, and $D_3$ are connected across a battery as shown in the figure. If the width of the depletion regions of $D_1$, $D_2$ and $D_3$ are $W_1$, $W_2$ and $W_3$, respectively, then the correct option is:
\fig{\begin{circuitikz}[scale=0.6,transform shape,font=\Large]
\draw (0,3.2) to[R,l=1\,k$\Omega$] (2.2,3.2) to[empty diode,l=$D_1$] (4.0,3.2) -- (4.4,3.2) -- (4.4,0.8);
\draw (0,3.2) -- (0,0.8);
\draw (0,2.0) to[R,l=1\,k$\Omega$] (2.2,2.0) to[empty diode,l=$D_2$] (3.7,2.0) node[circ]{}; \draw (4.1,2.0) node[circ]{} -- (4.4,2.0);
\draw (4.4,0.8) to[R,l_=1\,k$\Omega$] (2.2,0.8) to[empty diode,l_=$D_3$] (0,0.8);
\draw (0,2.0) -- (-0.8,2.0) -- (-0.8,-0.6) to[battery1,l_=5\,V] (5.2,-0.6) -- (5.2,2.0) -- (4.4,2.0);
\end{circuitikz}}\ntag{A}{335}{leph206}{13}{}{Class XII Part 2 \textperiodcentered{} Ch 14}}
{\opt{$W_1>W_2>W_3$}{$W_3 = W_1>W_2$}{$W_3>W_2>W_1$}{$W_2>W_1 = W_3$}}

\QQ{The following table presents the part of the electromagnetic spectrum and their corresponding major applications.\par\smallskip
{\small\begin{tabular}{@{}l@{\ }l@{\ }l@{\ }p{0.5\linewidth}@{}}
\multicolumn{2}{@{}l}{\textit{Part of the spectrum}} & \multicolumn{2}{l}{\textit{Application}}\\
P. & Microwave & I. & For purifying the water\\
Q. & UV rays & II. & For warming the food\\
R. & Gamma rays & III. & For AM and FM communication systems\\
S. & Radio wave & IV. & For treating the cancer cells
\end{tabular}}\par\smallskip
The correct option is:\ntag{D}{209}{leph108}{9}{ \textperiodcentered{} uses listed for each band}{Class XII Part 1 \textperiodcentered{} Ch 8}}
{\opt{P-I, Q-II, R-III, S-IV}{P-I, Q-IV, R-II, S-III}{P-II, Q-I, R-IV, S-III}{P-II, Q-IV, R-III, S-I}}

\QQ{Bob B of mass $m$ at rest is hanging vertically from the ceiling via a massless string of length $10\unit{m}$, as shown in the figure. Point mass A of mass $m$ travelling horizontally with speed $10\unit{m\,s^{-1}}$ hits bob B elastically. The bob B rises $h$ metre after the collision. Taking the acceleration due to gravity $g = 10\unit{m\,s^{-2}}$ and neglecting the size of the bob, the value of $h$ is:
\fig{\begin{tikzpicture}[scale=0.6]
\fill[pattern=north east lines] (-0.6,3) rectangle (0.6,3.25); \draw (-0.6,3) -- (0.6,3);
\draw (0,3) -- (0,0); \fill (0,0) circle (2.5pt) node[right,lbl]{B};
\fill (-1.6,0) circle (2pt) node[above,lbl]{A}; \draw[->] (-1.6,-0.3) -- (-0.9,-0.3) node[below,lbl]{$10\unit{m\,s^{-1}}$};
\end{tikzpicture}}\ntag{A}{84}{keph105}{14}{}{Class XI Part 1 \textperiodcentered{} Ch 5}}
{\opt{8}{7}{5}{2.5}}

\QQ{An ideal gas is made of polyatomic molecules. Each of the molecules has three translational, three rotational and $f$ number of vibrational modes. If the ratio of heat capacities $C_P/C_V$ of the gas is $8/7$, then the value of $f$ is:\ntag{D}{253}{keph205}{10}{ \textperiodcentered{} Cv=(3+f)R}{Class XI Part 2 \textperiodcentered{} Ch 12}}
{\opt{4}{3}{2}{1}}

%%BEGIN-ANALYSIS
\end{multicols}
\par\needspace{16\baselineskip}\vspace{14pt}
\begin{center}{\pgfornament[width=3.2cm,color=gilt]{87}}\par
{\large\color{navy}\scshape NEET 2026 \textperiodcentered{} Answer Key}\end{center}
{\small\setlength{\tabcolsep}{1pt}\renewcommand{\arraystretch}{1.5}
\begin{tabularx}{\linewidth}{@{}*{10}{>{\centering\arraybackslash}X}@{}}\toprule
\textbf{1}\,(4) & \textbf{2}\,(4) & \textbf{3}\,(4) & \textbf{4}\,(3) & \textbf{5}\,(1) & \textbf{6}\,(3) & \textbf{7}\,(1) & \textbf{8}\,(1) & \textbf{9}\,(4) & \textbf{10}\,(4)\\
\textbf{11}\,(1) & \textbf{12}\,(1) & \textbf{13}\,(3) & \textbf{14}\,(2) & \textbf{15}\,(2) & \textcolor{gray}{\textbf{16}\,(2)$^{\dagger}$} & \textbf{17}\,(3) & \textbf{18}\,(1) & \textbf{19}\,(3) & \textbf{20}\,(3)\\
\textbf{21}\,(4) & \textbf{22}\,(2) & \textbf{23}\,(3) & \textbf{24}\,(3) & \textbf{25}\,(4) & \textbf{26}\,--$^{*}$ & \textcolor{gray}{\textbf{27}\,(2)$^{\dagger}$} & \textbf{28}\,(2) & \textbf{29}\,(2) & \textbf{30}\,(1)\\
\textbf{31}\,(3) & \textbf{32}\,(2) & \textbf{33}\,(2) & \textbf{34}\,(3) & \textbf{35}\,(3) & \textbf{36}\,(1) & \textbf{37}\,(1) & \textbf{38}\,(2) & \textbf{39}\,(3) & \textbf{40}\,(4)\\
\textbf{41}\,(1) & \textbf{42}\,(3) & \textbf{43}\,(3) & \textbf{44}\,(3) & \textbf{45}\,(1) &  &  &  &  & \\\bottomrule\end{tabularx}}\par\smallskip
{\footnotesize Official NTA final key (Test Booklet Code 50); every answer was also solved and checked by the editor. (1)--(4) are the printed options. {\color{gray}Grey$^{\dagger}$}: deleted portion (answer given for completeness). --: question dropped by NTA (marks given to all). $^{*}$: see note.\par
\par\smallskip\textbf{Q26.} Dropped by NTA. The zero error is $-0.06$\,cm, so the actual length is 1.06\,cm, which is not among the options.
}\par
\clearpage
\begin{center}{\pgfornament[width=5cm,color=gilt]{87}}\par
{\fontsize{20}{24}\selectfont\color{navy}\scshape NEET 2026 \textperiodcentered{} Paper Analysis}\par
{\color{gilt}\pgfornament[width=5cm,color=gilt,symmetry=h]{87}}\end{center}
\vspace{4pt}
\begin{multicols}{2}
\textbf{\color{crimson}1. Direct vs Applied}\par\smallskip
{\footnotesize Checked question by question against the four 2026-27 NCERT books (XI Part 1 \& 2, XII Part 1 \& 2). Out of 45, 43 questions are in the present syllabus; percentages are of these.}\par\smallskip
\begin{tabular}{@{}lrr@{}}\toprule
Direct from NCERT & \textbf{16} & 37\%\\
\quad text, summary, table or figure & 15 &\\
\quad NCERT solved Example & 0 &\\
\quad NCERT Exercise & 1 &\\
Applied (concept in NCERT) & \textbf{27} & 63\%\\\midrule
Analysed & 43 &\\
Deleted portion (not analysed) & 2 &\\\midrule
\textbf{Total} & \textbf{45} &\\\bottomrule\end{tabular}\par\smallskip
{\footnotesize Direct: Q 2, 5, 7, 8, 12, 13, 14, 17, 24, 25, 35, 39, 40, 41, 43, 45.\\ From NCERT Exercises: Q 25 (Exercise 13.17).}\par\medskip
\textbf{\color{crimson}2. Difficulty level}\par\smallskip
\begin{tabular}{@{}lr@{}}\toprule
Easy & 16\\ Moderate & 27\\ Hard & 0\\\bottomrule\end{tabular}\par\smallskip
{\footnotesize The re-examination was moderate: most questions need one or two steps from a printed NCERT formula, and many test reading of graphs and figures. Class XI: 21 questions, Class XII: 22 (analysed questions only). (Difficulty is the editor's assessment.)}\par\medskip
\columnbreak
\textbf{\color{crimson}3. Chapter-wise distribution}\par\smallskip
{\footnotesize\begin{tabular}{@{}p{4.7cm}c@{\,}r@{\,}r@{\,}r@{}}\toprule
Chapter & Cl. & Qs & Dir & App\\\midrule
Units \& Measurement & XI & 3 & 0 & 3\\
Motion in a Straight Line & XI & 1 & 0 & 1\\
Laws of Motion & XI & 2 & 1 & 1\\
Work, Energy \& Power & XI & 2 & 0 & 2\\
Rotational Motion & XI & 1 & 0 & 1\\
Gravitation & XI & 2 & 2 & 0\\
Mech.\ Prop.\ of Fluids & XI & 2 & 1 & 1\\
Thermal Prop.\ of Matter & XI & 1 & 1 & 0\\
Thermodynamics & XI & 2 & 2 & 0\\
Kinetic Theory & XI & 2 & 2 & 0\\
Oscillations & XI & 2 & 1 & 1\\
Waves & XI & 1 & 0 & 1\\
Electrostatic Potential \& Capacitance & XII & 4 & 0 & 4\\
Current Electricity & XII & 1 & 0 & 1\\
Moving Charges \& Magnetism & XII & 2 & 0 & 2\\
Electromagnetic Induction & XII & 3 & 1 & 2\\
Alternating Current & XII & 1 & 0 & 1\\
Electromagnetic Waves & XII & 2 & 1 & 1\\
Ray Optics & XII & 2 & 0 & 2\\
Dual Nature & XII & 3 & 1 & 2\\
Atoms & XII & 2 & 2 & 0\\
Nuclei & XII & 1 & 1 & 0\\
Semiconductor Electronics & XII & 1 & 0 & 1\\
\bottomrule\end{tabular}}\par\medskip
\textbf{\color{crimson}4. Teaching takeaways}\par\smallskip
{\footnotesize\begin{itemize}[leftmargin=1em,nosep]
\item Thermodynamics and kinetic theory gave quick direct marks: adiabatic work, cyclic process, mean free path and $C_v=(3+f)R$ are printed formulas.
\item Exercise 13.17 (floating cork in SHM) was asked almost unchanged; NCERT exercises must be solved fully.
\item Two practical-skill questions (index correction, vernier zero error) need the Lab Manual; the vernier question was dropped by NTA.
\item Graph-reading questions (photocurrent vs time, terminal velocity vs $\sigma/\rho$, $N(\theta)$ in Geiger--Marsden) come straight from NCERT figures.
\end{itemize}}\par\medskip
\textbf{\color{crimson}5. Deleted portions}\par\smallskip
\textbf{2} of the 45 questions are on topics removed from the 2026-27 NCERT (rationalised syllabus). They are marked \textsc{Deleted Portion} and are not analysed: Zener diode (Q 16); parallel axis theorem (Q 27).\par
\end{multicols}
%%END-ANALYSIS


